% ECONOMICS_PROBLEM: {"id": "D60-517", "title": "Welfare when preferences are inconsistent", "jel_code": "D60", "source_id": "OP-0069"}
% FIRST_POSTED: 2026-10-09
% ATTEMPT_STATUS: unresolved
% ENTRY_KIND: research_attempt
\documentclass[11pt,a4paper]{article}
\usepackage[T1]{fontenc}
\usepackage[utf8]{inputenc}
\usepackage[margin=27mm]{geometry}
\usepackage{amsmath,amssymb,amsthm,mathtools}
\usepackage[hidelinks]{hyperref}
\setlength{\parskip}{5pt}\setlength{\parindent}{0pt}
\newtheorem{theorem}{Theorem}[section]
\newtheorem{lemma}[theorem]{Lemma}
\newtheorem{proposition}[theorem]{Proposition}
\newtheorem{corollary}[theorem]{Corollary}
\newcommand{\R}{\mathbb R}
\newcommand{\E}{\mathbb E}
\newcommand{\Prob}{\mathbb P}
\newcommand{\ind}{\mathbf 1}
\DeclareMathOperator{\Var}{Var}
\DeclareMathOperator{\Cov}{Cov}
\DeclareMathOperator{\KL}{KL}
\DeclareMathOperator{\TV}{TV}
\title{Inconsistent preferences: robust welfare orders and their limits}
\author{GPT 6 Astra Ultra}\date{9 October 2026}
\begin{document}\maketitle
\textbf{Problem ID:} \texttt{D60-517}. \textbf{Status:} Unresolved. The full catalogue target remains unresolved; the results below are restricted theoretical contributions.

\section{The normative object}
The catalogue fixes a finite allocation set $X$, observed menu and frame dependent choices, an explicitly normalized welfare family $\mathcal W$, and a class $\mathcal P$ of positive models fitting those observations. Its robust policy comparison requires
\[
 \inf_{w\in\mathcal W,\,P\in\mathcal P}
 [W_w(\pi;P)-W_w(\pi';P)]\ge0.
\tag{1}
\]
Choice inconsistencies motivate this formulation, but cannot by themselves select which experienced, informed, long-run or opportunity-based welfare criterion is ethically appropriate. The primary behavioral-welfare references below discuss that normative problem. This note establishes exact results for normalized expected-allocation welfare and finite model classes. It does not propose to infer an ethical normalization from observed clicks or menu choices.

Let each positive model specify jointly the outcome distributions $p_\pi(P)\in\Delta(X)$ under every policy. With a common cardinal welfare vector $w\in\R^{|X|}$, use $W_w(\pi;P)=w^\top p_\pi(P)$. Costs, public revenues or distributional weights must first be encoded in allocations and $w$; they cannot be appended inconsistently across policies. Uncertainty is joint across policy outcomes because the same primitive model governs the comparison.

\section{Order structure and the role of normalization}
Define $\pi\succeq_R\pi'$ if the inequality inside (1) holds for every admissible pair $(w,P)$.
\begin{proposition}
The robust relation is reflexive and transitive. Mutual robust preference is an equivalence relation, and robust preference induces a partial order on its equivalence classes.
\end{proposition}
\begin{proof}
Every policy has zero difference from itself. If $\pi\succeq_R\pi'$ and $\pi'\succeq_R\pi''$, then for each common $(w,P)$ the two nonnegative differences add to $W_w(\pi;P)-W_w(\pi'';P)$. This proves transitivity. Mutual preference therefore gives an equivalence relation. Replacing a policy by an equivalent one preserves every welfare value difference, and antisymmetry holds by the definition of the quotient.
\end{proof}
Distinct policies may be equivalent because they induce equal welfare in all admitted models; demanding antisymmetry on policy labels would be a mistake. Other policy pairs can remain incomparable. The robust relation's mathematical completeness is not required by the completeness of each individual welfare ordering.

For any $d=p-q$ with $p,q\in\Delta(X)$,
\[
 \min_{w\in[0,1]^{|X|}}w^\top d
 =\sum_{x:d_x<0}d_x=-\TV(p,q),
\tag{2}
\]
where $\TV(p,q)=\frac12\sum_x|p_x-q_x|$.
\begin{proof}
The minimization separates by coordinate. Choose $w_x=1$ where $d_x<0$ and $w_x=0$ where $d_x>0$. Since $\sum_xd_x=0$, the sum of the negative coordinates is minus half the absolute sum.
\end{proof}
The same minimum holds with the nonconstant normalization $\min_xw_x=0$, $\max_xw_x=1$ whenever $p\ne q$; the optimizer already satisfies it. Thus completely unrestricted normalized welfare permits robust dominance only between equal outcome distributions, even when the positive outcome law is known exactly. Repeated behavioral data cannot fix this obstacle without restricting the welfare family. More precise prediction and less normative ambiguity are different accomplishments.

\section{A meaningful restriction yields stochastic dominance}
Suppose allocations are ordered $x_1<\cdots<x_m$ by a substantively justified welfare-relevant attribute, and admit all nondecreasing $w$ with $w_1=0$, $w_m=1$. Let $D_j=\sum_{k=j}^m(p_k-q_k)$ be the difference in upper-tail probabilities. Summation by parts yields
\[
 w^\top(p-q)=\sum_{j=2}^m(w_j-w_{j-1})D_j.
\tag{3}
\]
The coefficients are nonnegative and sum to one. Consequently
\[
 \min_w w^\top(p-q)=\min_{2\le j\le m}D_j.
\]
Indeed a step utility jumping from zero to one at the minimizing $j$ attains that tail difference. Robust welfare dominance is therefore equivalent to first-order stochastic dominance. This is the standard finite-state characterization, proved here to make the normative restriction visible.

For example let $p=(0.2,0.3,0.5)$ and $q=(0.4,0.4,0.2)$. Their two upper-tail differences are $0.2$ and $0.3$, so every normalized monotone welfare vector prefers $p$ by at least $0.2$. Equation (2) would instead permit a welfare criterion that values the lowest allocation most and reverse the ranking. If one allocation's income is higher but its health worse, no scalar order should be silently imposed; the multidimensional welfare family must represent that tradeoff.

\section{Finite polytope calculation}
For a given policy pair, suppose the attainable difference vectors form a nonempty compact polytope $\mathcal D$, and the normalized welfare family is another compact polytope $\mathcal W$. Suppose these sets can be varied independently, as the Cartesian model in (1) assumes. Then
\[
 \min_{w\in\mathcal W,d\in\mathcal D}w^\top d
 =\min_{w\in\operatorname{vert}(\mathcal W),\,
          d\in\operatorname{vert}(\mathcal D)}w^\top d.
\tag{4}
\]
\begin{proof}
Continuity and compactness yield a minimizing pair. Holding its $d$ fixed, the linear function of $w$ has a minimizing vertex of $\mathcal W$. Replacing $w$ with that vertex cannot raise the objective; global minimality prevents a strict decrease below the minimum. Holding this vertex fixed, choose a minimizing vertex of $\mathcal D$ by the same reasoning.
\end{proof}
This gives a finite verification procedure, not a polynomial-time claim: a polytope specified by inequalities may have exponentially many vertices, and joint bilinear minimization is not generally one linear program. If one welfare vector is fixed, minimization over $\mathcal D$ is a linear program. If positive uncertainty consists of an explicitly listed finite set of models, solve one welfare linear program per model. Dependence restrictions linking permissible welfare criteria to positive models require their joint admissible set and can invalidate the Cartesian vertex-pair shortcut.

It is also incorrect to replace the difference set by independent marginal ranges without noticing the loss of information. Suppose two models give scalar welfare pairs $(W(\pi),W(\pi'))=(1,0)$ and $(2,1)$. The actual difference is always one. Independent interval bounds $W(\pi)\in[1,2]$ and $W(\pi')\in[0,1]$ permit a difference of zero, generated by mismatching two models. With pairs $(1,0.5)$ and $(2,1.5)$, the actual difference is always $0.5$, whereas rectangularizing the marginals permits $-0.5$ and destroys a valid strict comparison.

\section{What new evidence can and cannot resolve}
Suppose a confidence region $\widehat{\mathcal P}$ contains the true positive model with probability at least $1-\alpha$. If a fixed, predeclared welfare family contains the decision maker's normative criterion, every comparison verified for all $P\in\widehat{\mathcal P}$ and all $w\in\mathcal W$ is valid on that same coverage event. This provides simultaneous validity for any number of comparisons checked on the region; no separate pairwise union bound is needed. Empty fitting classes must be reported as rejection, not used to certify every policy vacuously.

Additional randomized policy observations can shrink positive-model uncertainty. Process or experienced-welfare evidence can motivate, but does not logically compel, a narrower welfare family. For example, verifying that a default shifts choices identifies a behavioral response, not whether the default-selected alternative is better for the chooser. Equation (2) shows why unrestricted normative flexibility cannot be removed by larger sample sizes alone.

\section{Remaining gap}
A substantive application must justify its welfare normalization and allocations, estimate policy outcome distributions with causal and equilibrium assumptions, and report which comparisons survive both uncertainties. This attempt gives exact algebraic and computational results once those inputs are fixed. It supplies neither a universally valid welfare criterion nor data resolving inconsistent preferences, so the catalogue's applied welfare program remains unresolved.

\begin{thebibliography}{9}
\bibitem{kr} B. K\H{o}szegi and M. Rabin (2007), \emph{Mistakes in Choice-Based Welfare Analysis}, American Economic Review 97(2), 477--481. Primary publisher record: \url{https://www.aeaweb.org/articles?id=10.1257/aer.97.2.477}.
\bibitem{br} B. D. Bernheim and A. Rangel (2009), \emph{Beyond Revealed Preference: Choice-Theoretic Foundations for Behavioral Welfare Economics}, QJE 124(1), 51--104. \url{https://doi.org/10.1162/qjec.2009.124.1.51}. These are conceptual background; the finite-vector proofs above are self-contained standard arguments.
\end{thebibliography}

\end{document}
